\section{Generated Scanners} 

	It turns out that using handwritten lexers is not the most powerful nor practical way to do this. Remember that all scanning does is pattern matching using \hyperref[th-def:regular_expression]{regular expressions}. In modern days, we have software that can take in whatever set of regular expressions you give it, and it will convert it to an \hyperref[th-thm:regular_equivalence]{equivalent} \hyperref[th-def:dfa]{DFA}. Therefore, to detect tokens, we run it through a DFA and see the accepting state it lands on. 

  Again, with multiple matches, the problem is manifested with multiple dead states. If it is in a dead state, then output the most recent\footnote{due to principle of maximal munch} accepting state. 

  \begin{example}[Ambiguities in Tokenizing]
    Say that we have a language with the following tokens. 
    \begin{enumerate}
      \item \texttt{<x>}: \texttt{x}. 
      \item \texttt{<nxy1>}: \texttt{(x*)y} 
      \item \texttt{<z>}: \texttt{z}
    \end{enumerate}
    This leads to the following DFA. 

    \begin{figure}[H]
      \centering 
      \begin{tikzpicture}[automata]
        % States
        \node[state, initial, initial text=] (q0) {};
        \node[state, accepting, right=of q0] (qx) {\texttt{<x>}};
        \node[state, accepting, below=of q0] (qz) {\texttt{<z>}};
        \node[state, accepting, below=of qx] (qnxy) {\texttt{<nxyl>}};
        \node[state, right=of qx] (qtrap) {};

        % Transitions
        \path (q0) edge node {\texttt{x}} (qx)
              (q0) edge node {\texttt{z}} (qz)
              (q0) edge node {\texttt{y}} (qnxy)
              (qx) edge node {\texttt{x}} (qtrap)
              (qx) edge node {\texttt{y}} (qnxy)
              (qtrap) edge [loop above] node {\texttt{x}} (qtrap)
              (qtrap) edge node[swap] {\texttt{y}} (qnxy);
      \end{tikzpicture}
      \caption{} 
      \label{fig:dfa3choice}
    \end{figure}

    Then, to parse ``\texttt{xxxzxxxy},'' we have 3 choices, which leads to ambiguity. 
    \begin{enumerate}
      \item \texttt{<x> <x> <x> <z> <nxy1>}
      \item \texttt{<x> <x> <x> <z> <x> <nxy1>}
      \item \texttt{<x> <x> <x> <z> <x> <x> <nxy1>}
    \end{enumerate}
    Again, we can use maximal munch.  
  \end{example}

  Lexers can have multiple DFAs, and depending on context (through start state), you traverse different DFAs (e.g. code vs comments). 

